% Generated draft: requires mathematical and source review.
\documentclass[11pt]{article}
\usepackage[margin=25mm]{geometry}
\usepackage{fontspec}
\setmainfont{Latin Modern Roman}
\usepackage{amsmath,amssymb,mathtools,unicode-math}
\setmathfont{Latin Modern Math}
\usepackage{enumitem,needspace,xurl,hyperref}
\setlength{\parindent}{0pt}
\setlength{\parskip}{5pt}
\setlength{\emergencystretch}{3em}
\newcommand{\N}{\mathbb{N}}
\newcommand{\Z}{\mathbb{Z}}
\newcommand{\Q}{\mathbb{Q}}
\newcommand{\R}{\mathbb{R}}
\newcommand{\C}{\mathbb{C}}
\newcommand{\F}{\mathbb{F}}
\newcommand{\A}{\mathbb{A}}
\newcommand{\PP}{\mathbb{P}}
\newcommand{\E}{\mathbb{E}}
\newcommand{\eps}{\varepsilon}
\newcommand{\dd}{\,\mathrm{d}}
\begin{document}
\section*{UnsolvedMath ID: 2306105}
\textbf{Source problem code:} catalogue entry
\subsection{Definitions and mathematical statement}
In the usual harmonic-univalent setting, write $f=h+\overline g$ on $\mathbb D$, with $h,g$ holomorphic. Sense preservation means $|h'|^2-|g'|^2>0$. The class $S_H$ conventionally consists of normalized sense-preserving univalent harmonic maps; the source does not give its exact normalization, including whether $g'(0)=0$ is required. Let
\[
K_H=\{f\in S_H:f(\mathbb D)\text{ is convex}\}.
\]
For holomorphic series $u(z)=\sum u_nz^n$, $v(z)=\sum v_nz^n$, their Hadamard convolution is $(u*v)(z)=\sum u_nv_nz^n$. A harmonic coefficient multiplier acts on the analytic and coanalytic coefficients by prescribed sequences. The source's symbol $\phi^*$ does not specify whether those sequences arise from one analytic kernel or from a pair of harmonic kernels, so that convention must be fixed.

\textbf{Multiplier problem (T. Sheil-Small).} Characterize all convolution operators $\phi^*$ for which
\[
\phi^*(f)\in K_H\qquad\text{for every }f\in K_H.
\]
\textbf{Radius problem.} Determine exactly the set
\[
\mathcal R=\{r\in(0,1):f(r\mathbb D)\text{ is convex for every }f\in K_H\}.
\]
This is the literal universal reading of asking when $z\mapsto f(rz)$ is convex in the disk. It concerns convexity of its image, not convexity of the function as a real-valued scalar function. If normalization is needed for class membership, $f(rz)/r$ preserves the same convexity property and the usual derivative-one normalization; the source's map itself is $f(rz)$.

The source reports the restriction $r\leq\sqrt2-1$, attributed to Clunie and Sheil-Small. This is retained in its stated upper-bound direction, not silently changed to a sufficient-radius theorem. The supplied summary reports no definitive resolution located as of its 2018-based assessment.
\subsection{Short English statement}
Find all coefficient-convolution operators preserving convex harmonic univalent maps, and determine the universal radii for which every such map remains convex on the smaller disk. Kernel and normalization conventions are missing.
\subsection{Sources}
\begin{itemize}
\item Hayman, \emph{Research Problems in Function Theory} (2018), Problem 6.105, attributed to T. Sheil-Small, \url{https://arxiv.org/abs/1809.07200}; source-TeX extraction, accessed 2026-07-29. The 2018 update was retained for status triage.
\item Clunie and Sheil-Small, the stated $r\leq\sqrt2-1$ radius restriction; the extracted reference has no title, date, or theorem location.
\item Local registry \texttt{data/unsolvedmath/problems.json}, permanent ID $2306105$, source code catalogue entry, \emph{Research Problems in Function Theory --- Problem 6.105}. The record labels this open and OPEN-TRIAGE, with NEEDS\_REVIEW source-status and rights qualifications.
\end{itemize}
\end{document}
